\noindent\textbf{Resultado.} $y\equiv0$ y $y=\pm[x/2-1/20+Ce^{-10x}]^{-1/2}$ en cada intervalo donde el corchete sea positivo.

\noindent\textit{Desarrollo.} La función cero es solución. Para una rama no nula, $z=y^{-2}$ da $z'=-2y^{-3}y'$ y transforma la ecuación en $z'+10z=5x$. El factor integrante $e^{10x}$ conduce a $z=x/2-1/20+Ce^{-10x}$. Puesto que $z=y^{-2}>0$, sólo valen los intervalos donde esta expresión es positiva; en ellos se toma un signo fijo para $y$. La sustitución inversa verifica la EDO. El campo original es suave en toda la recta del plano, así que la unicidad impide que una rama no nula se una a la solución cero en un punto interior. Los ceros del corchete son extremos con divergencia de $|y|$.
